\chapter{REVISIÓN DE LITERATURA}

\section{Antecedentes}

\subsection{Antecedentes internacionales}
\begin{enumerate}

\item Abou Elasaad et al. (2025), en su artículo AegisGuard: A Multi-Stage Hybrid Intrusion Detection System with Optimized Feature Selection for Industrial IoT Security, propusieron un estudio experimental orientado a fortalecer ecosistemas IIoT mediante arquitecturas de análisis secuencial. La metodología implementó AegisGuard, un modelo modular de múltiples etapas que filtra flujos maliciosos aplicando selección dinámica de características. Los resultados evidenciaron empíricamente que este enfoque estratificado supera estadísticamente a los sistemas de etapa única en la neutralización asertiva de ciberataques estructuralmente complejos. Este trabajo respalda concluyentemente el paradigma de construir arquitecturas híbridas y multicapa; sin embargo, no atiende primariamente el sesgo algorítmico generado por clases minoritarias subrepresentadas en la capa de entrada, variable crítica que esta tesis subsanará empleando el algoritmo Borderline-SMOTE.

\item Adewole et al. (2025), en Intrusion Detection Framework for Internet of Things with Rule Induction for Model Explanation, ejecutaron un diseño experimental para integrar capacidades de explicabilidad en sistemas de ciberseguridad IoT. Metodológicamente, utilizaron los conjuntos de datos CIC-IDS2017 y CICIoT2023 para evaluar el rendimiento de algoritmos de Gradient Boosting, incluyendo XGBoost, LightGBM y CatBoost. Los resultados demostraron la superioridad técnica de XGBoost, el cual alcanzó un Accuracy de 99.91% en el primer dataset, 98.54% en el CICIoT2023 y un AUC-ROC global de 93.06%. Este trabajo justifica rigurosamente la elección de XGBoost como algoritmo base competente; no obstante, no implementa técnicas de rebalanceo avanzado, requerimiento metodológico indispensable que la presente investigación abordará.

\item Ahmed et al. (2024), en An optimized ensemble model with advanced feature selection for network intrusion detection, realizaron un estudio de nivel explicativo con el objetivo de optimizar la identificación de amenazas mediante la hibridación de algoritmos predictivos. La metodología consistió en integrar los modelos XGBoost y LightGBM bajo una arquitectura de ensamble tipo Stacking, evaluando el impacto de la remoción de atributos irrelevantes en la estabilidad sistémica. Los resultados confirmaron que la mitigación del ruido de entrada acelera significativamente el entrenamiento e incrementa la precisión global del clasificador. La relevancia de esta investigación radica en validar la efectividad estructural del modelo Stacking seleccionado; sin embargo, carece de un enfoque orientado específicamente a métricas de ataques minoritarios, dimensión que será el núcleo de evaluación del modelo propuesto en esta tesis.

\item Al-Emari et al. (2024), en su artículo Employing Mutual Information Feature Selection and LightGBM for Intrusion Detection in IoT, desarrollaron una investigación aplicada cuyo objetivo fue optimizar los mecanismos de protección de redes mediante la selección heurística de características. Metodológicamente, emplearon el algoritmo LightGBM evaluado sobre el conjunto de datos IoTID20, integrando el análisis de Información Mutua y optimización por enjambre de partículas. Los resultados revelaron que esta configuración supera a las técnicas convencionales en las métricas de exactitud, sensibilidad y F1-Score, reduciendo simultáneamente la tasa de falsos positivos. Este estudio demuestra fehacientemente la eficacia del clasificador LightGBM para entornos IoT; sin embargo, al utilizar un dataset genérico y optimizaciones de alto costo, deja un vacío técnico que la presente tesis cubrirá mediante el dataset CICIoT2023.

\item Almotairi et al. (2024), en Enhancing intrusion detection in IoT networks using machine learning-based feature selection and ensemble models, desarrollaron una investigación aplicada sobre la gestión algorítmica ante la alta heterogeneidad del tráfico en sistemas IoT. Su protocolo metodológico consistió en acoplar modelos de ensamble clásicos, como Random Forest y Gradient Boosting, con rutinas selectivas de atributos. Los hallazgos revelaron que una discriminación óptima de variables permite a las arquitecturas mantener tasas de exactitud excepcionales garantizando baja demanda de hardware de procesamiento. Este trabajo consolida la viabilidad de utilizar algoritmos basados en árboles de decisión en ecosistemas restrictivos; no obstante, carece de evaluación sobre conjuntos de ultra-alta dimensionalidad modernos, limitación que se superará al integrar el modelo con los 46 millones de flujos del CICIoT2023.

\item Alrayes et al. (2024), en Intrusion Detection in IoT Systems Using Denoising Autoencoder, desarrollaron una investigación aplicada enfocada en la seguridad para ecosistemas IoT sujetos a alta variabilidad y recursos limitados. Utilizaron Autoencoders de Eliminación de Ruido (DAE) basados en aprendizaje no supervisado para la extracción de características, validando el sistema con los conjuntos de datos CICIDS2017 y NSL-KDD. Los resultados exhibieron precisiones sobresalientes del 99.991% y 99.4% respectivamente, demostrando alta fiabilidad en procesamiento de eventos en tiempo real. Este trabajo resalta la relevancia de filtrar atributos para entornos de bajo consumo; sin embargo, su fundamentación en redes neuronales valida la propuesta de la presente tesis de integrar enfoques Stacking más ligeros para garantizar latencias de inferencia de rango inferior.

\item Berguiga et al. (2025), en HIDS-IoMT: A Deep Learning-Based Intelligent Intrusion Detection System for the Internet of Medical Things, plantearon un estudio experimental orientado a proteger redes de dispositivos médicos operando sobre arquitecturas computacionales de niebla. Metodológicamente combinaron redes convolucionales (CNN) y recurrentes (LSTM) para clasificar anomalías, ejecutando evaluaciones sobre los conjuntos de datos IoTID20 y Edge-IIoTset. Los hallazgos mostraron métricas críticas de Accuracy (99.92%) y Recall (99.99%) ante ataques disruptivos como denegación de servicio. Este antecedente subraya la robustez de los modelos híbridos para infraestructuras IoT de misión crítica; sin embargo, se enfoca primordialmente en ataques volumétricos mayores, dejando un área de oportunidad técnica en la identificación estricta de ataques minoritarios encubiertos.

\item Diwan et al. (2022), en su investigación Feature Entropy Estimation (FEE) for Malicious IoT Traffic and Detection Using Machine Learning, condujeron un estudio cuantitativo orientado a mejorar la identificación de actividades anómalas en volúmenes masivos de tráfico de red. La metodología propuesta consistió en utilizar la estimación de entropía para la ingeniería de características, evaluando algoritmos como XGBoost y Random Forest sobre el conjunto de datos Bot-IoT. Los hallazgos confirmaron que el uso de la entropía para discriminar atributos permite a los modelos alcanzar una precisión superior al 97% en la clasificación de amenazas. Este antecedente resulta fundamental para justificar el uso de la entropía de Shannon; no obstante, su enfoque no aborda la mitigación del desbalance extremo de clases, requerimiento que será subsanado.

\item Ferrag et al. (2022), en su artículo Edge-IIoTset: A New Comprehensive Realistic Cyber Security Dataset of IoT and IIoT Applications for Centralized and Federated Learning, ejecutaron un diseño experimental con el objetivo de proponer y validar un nuevo conjunto de datos representativo de amenazas modernas. Metodológicamente, generaron perfiles de tráfico desde múltiples dispositivos físicos, extrayendo características correlacionadas para evaluar el rendimiento de algoritmos tradicionales y enfoques profundos. Los resultados evidenciaron que los modelos predictivos estándar presentan severas limitaciones de exactitud cuando se enfrentan a escenarios realistas de intrusión. La relevancia de este estudio estriba en demostrar la insufficiency de los clasificadores individuales ante la complejidad del tráfico, justificando la necesidad de implementar metamodelos robustos.

\item Gutierrez et al. (2023), en Enhancing Intrusion Detection in IoT Communications Through ML Model Generalization With a New Dataset (IDSAI), ejecutaron una investigación aplicada destinada a cuantificar la capacidad de generalización en modelos de aprendizaje supervisado. La metodología introdujo el dataset normado y balanceado IDSAI para entrenar Random Forest, Gradient Boosting y XGBoost, validándolos posteriormente sobre el dataset externo Bot-IoT. Los hallazgos confirmaron a XGBoost como el modelo de mayor eficacia interna (94.97%), preservando un desempeño notable frente a datos externos no contemplados. Este artículo cimenta la elección de XGBoost por su alta resiliencia predictiva; no obstante, al utilizar un entorno pre-balanceado, omite el tratamiento del desbalance de tráfico real que abordará la presente tesis.

\item Kouassi et al. (2025a), en Optimized Intrusion Detection in the IoT Through Statistical Selection and Classification with CatBoost and SNN, condujeron un estudio cuantitativo orientado a mejorar la detección de anomalías abordando directamente el problema del desbalance de clases. En su marco metodológico, aplicaron técnicas de remuestreo como SMOTE y RUS de forma conjunta con regresión LASSO, validando la arquitectura CatBoost en el dataset IoTID20. Los resultados indican que la configuración híbrida maximiza el rendimiento, logrando una precisión de 82.43% y un F1-Score de 85.08%. El estudio demuestra empíricamente la imperatividad del rebalanceo previo al entrenamiento; sin embargo, evalúa arquitecturas alternativas en un dataset de menor dimensionalidad que el CICIoT2023, brecha técnica que será suplida por el presente proyecto.

\item Kouassi et al. (2025b), en Top-K Feature Selection for IoT Intrusion Detection: Contributions of XGBoost, LightGBM, and Random Forest, desarrollaron una investigación aplicada orientada a reducir la carga computacional algorítmica en entornos IoT. Emplearon el dataset CICIoMT2024 aplicando una selección estadística Top-K sobre clasificadores supervisados, restringiendo las variables a las 10 y 15 más relevantes. Los hallazgos demostraron que el algoritmo Random Forest con Top-10 alcanzó un Accuracy de 91.7% y un F1-Score de 93%, reduciendo el tiempo de procesamiento en un 35%. Este estudio corrobora la pertinencia de la selección de características; no obstante, su enfoque en clasificadores individuales deja un vacío respecto a la capacidad predictiva de las arquitecturas Stacking.

\item Lin et al. (2021), en el documento Ensemble Learning for Threat Classification in Network Intrusion Detection on a Security Monitoring System for Renewable Energy, desarrollaron una investigación aplicada para categorizar amenazas informáticas mediante esquemas de aprendizaje por ensamble. En su metodología, contrastaron el rendimiento de algoritmos basados en bosques aleatorios frente a modelos de Gradient Boosting como XGBoost, procesando registros de tráfico de red para identificar anomalías. Los resultados comprobaron que los enfoques de ensamble garantizan una alta precisión y sensibilidad en la detección. Este antecedente consolida la viabilidad técnica de integrar algoritmos como XGBoost para tareas de clasificación crítica; sin embargo, presenta una oportunidad técnica para extender y optimizar esta arquitectura hacia ecosistemas IoT con la técnica SMOTE.

\item Musthafa et al. (2024), en Optimizing IoT Intrusion Detection Using Balanced Class Distribution, Feature Selection, and Ensemble Machine Learning Techniques, realizaron un estudio cuantitativo orientado a optimizar la fiabilidad de los sistemas de detección. La metodología consistió en implementar arquitecturas Stacking con LSTM, aplicando un balanceo de clases y selección de características sobre los conjuntos UNSW-NB15 y NSL-KDD. Los hallazgos indican que la configuración Stacking con filtrado estadístico alcanza una precisión de hasta 99.77%, mitigando significativamente el sobreajuste. Este artículo fundamenta empíricamente que la sinergia entre rebalanceo y Stacking maximiza el rendimiento predictivo; no obstante, la utilización de Deep Learning conlleva una alta carga computacional que la presente tesis evitará mediante algoritmos ligeros.

\item Nigar y Mustafa (2025), en Enhanced Intrusion Detection via Hybrid Data Resampling and Feature Optimization, elaboraron un diseño experimental combinando técnicas de optimización de atributos y sobremuestreo para estabilizar algoritmos predictivos. La metodología integró el algoritmo SMOTE para mitigar el desbalance previo al entrenamiento de ensambles modernos, particularmente XGBoost y LightGBM. Los resultados arrojaron que este preprocesamiento conjunto reduce significativamente el ratio de falsos positivos y proporciona alta estabilidad analítica frente a vectores de ataque de día cero. Este estudio proporciona un respaldo metodológico robusto; sin embargo, no contempla metadatos de incertidumbre lógica (entropía) para asistir la predicción, confirmando el vacío técnico que resolverá la presente propuesta de Stacking.

\item Popoola et al. (2025), en Multi-Stage Deep Learning for Intrusion Detection in Industrial Internet of Things, propusieron una investigación aplicada para resolver la identificación de ataques en entornos IIoT altamente desbalanceados. Su metodología incorporó un enfoque de aprendizaje profundo multi-etapa que utiliza la métrica de entropía para discriminar la incertidumbre predictiva, evaluado en los datasets X-IIOTID y WUSTL-IIoT. Los resultados evidenciaron métricas de precisión superiores al 99%, constatando que la cuantificación de la incertidumbre facilita la detección de intrusiones complejas y poco frecuentes. Este antecedente fundamenta teóricamente el uso de la entropía de Shannon en el metamodelo propuesto; no obstante, motiva el diseño de una solución más ligera basada en Stacking.

\item Sama et al. (2025), en su artículo Cutting-Edge Intrusion Detection in IoT Networks: A Focus on Ensemble Models, ejecutaron un estudio cuantitativo cuyo objetivo fue optimizar la detección de intrusiones mediante modelos de ensamble. Metodológicamente, emplearon datasets masivos de tráfico de red para evaluar la arquitectura híbrida Stacking, combinando XGBoost y LightGBM junto con métricas de incertidumbre (entropía). Los resultados evidenciaron que este modelo supera estadísticamente a los clasificadores individuales en las métricas de Accuracy y Recall, minimizando simultáneamente las falsas alarmas. Este antecedente sustenta empíricamente la superioridad del ensamble propuesto; sin embargo, no aborda el desbalance extremo de datos, justificando la inclusión de la técnica Borderline-SMOTE en la presente investigación.

\item Sayegh et al. (2024), en Enhanced Intrusion Detection with LSTM-Based Model, Feature Selection, and SMOTE for Imbalanced Data, realizaron un estudio cuantitativo orientado a elevar las tasas de detección mitigando el sesgo algorítmico. La metodología integró el remuestreo SMOTE con selección de características RFE, evaluando el desempeño en datasets heterogéneos (CIC-IDS2017, NSL-KDD, UNSW-NB15). Los resultados confirmaron que la aplicación de SMOTE incrementa drásticamente la capacidad de detección, alcanzando precisiones pico de 99.34%. Este estudio sustenta la eficacia innegable del sobremuestreo sintético; no obstante, emplea una variante estándar, lo que justifica la implementación metodológica de Borderline-SMOTE para focalizar la generación de datos.

\item Sepúlveda-Fontaine y Amigó (2024), en el estudio teórico Applications of Entropy in Data Analysis and Machine Learning: A Review, condujeron una investigación documental con el objetivo de sistematizar el uso de medidas entrópicas en la optimización algorítmica. Metodológicamente, revisaron la aplicación de la entropía de Shannon como un funcional axiomático para caracterizar distribuciones de masa de probabilidad en el contexto de la minería de datos. Sus conclusiones determinaron que la entropía resulta excepcionalmente adecuada para evaluar la importancia de las características y medir el error en modelos predictivos. Este antecedente provee el sustento teórico y matemático necesario para justificar la integración de la Entropía de Shannon como parámetro de evaluación en el metamodelo propuesto.

\item Shanmugam et al. (2025), en Addressing Class Imbalance in Intrusion Detection: A Comprehensive Evaluation of Machine Learning Approaches, condujeron un estudio de diseño explicativo enfocado en evaluar rigurosamente el impacto de las estrategias de preprocesamiento. Emplearon métodos de rebalanceo algorítmico, como SMOTE y ADASYN, contrastados estadísticamente frente a diversos clasificadores estándar. Los resultados determinaron que el remuestreo incrementa drásticamente la tasa de detección de anomalías infrecuentes sin degradar el rendimiento métrico en clases mayoritarias. Este antecedente resulta teóricamente vinculante con el presente proyecto al validar la necesidad ineludible del preprocesamiento para ataques minoritarios.

\item Wu et al. (2025), en MARNet: An Efficient Two-Stage Intrusion Detection Model Based on Deep Learning, diseñaron un estudio experimental para optimizar la identificación de vectores de ataque infrecuentes en redes informáticas. Emplearon los datasets UNSW-NB15 y NSL-KDD implementando la función de pérdida EQL v2 sobre redes neuronales residuales para mitigar el efecto negativo del desbalance de clases. La arquitectura propuesta alcanzó una precisión multiclase del 90.55% y una tasa de detección del 99.51%. Este artículo valida la criticidad de gestionar el desbalance para detectar intrusiones minoritarias; sin embargo, al emplear estructuras de Deep Learning, presenta limitaciones de viabilidad operativa que se abordarán empleando algoritmos de ensamble.

\item Zwane et al. (2019), en Ensemble Learning Approach for Flow-based Intrusion Detection System, ejecutaron un estudio experimental destined to improve the analysis of network flows through compound predictive methods. In their methodological design, they evaluated probabilistic and non-probabilistic ensemble techniques based on decision trees to classify intrusions. The results established that ensemble methods systematically surpass simple algorithms in generalization capability over anomalous traffic. Although this study establishes the fundamental basis for the use of combined architectures, its analysis lacks specific mechanisms for treating low-frequency attacks, a methodological gap that will be resolved through the proposed hybrid approach.

\end{enumerate}

\subsection{Antecedentes nacionales}
\begin{enumerate}

\item Raico Gallardo (2025), en Desarrollo de un modelo de clasificación para la detección de anomalías en redes IoT utilizando el dataset CICIoT2023, desarrolló un estudio aplicado de enfoque cuantitativo con el fin de identificar patrones maliciosos en entornos de Internet de las Cosas. En su diseño metodológico, procesó una muestra de 7,829,178 registros del dataset CICIoT2023, evaluando algoritmos como Regresión Logística, Árboles de Decisión y Random Forest mediante la herramienta Scikit-learn. Los resultados determinaron que el modelo Random Forest alcanzó la mayor efectividad con un 98.63% de exactitud, consolidándose como una técnica robusta para la clasificación de tráfico. Si bien esta investigación valida la potencia del dataset que se utilizará en el presente estudio, se identifica un vacío en el tratamiento de ataques de baja frecuencia, una oportunidad que será abordada mediante la integración de Stacking y técnicas de sobremuestreo sintético.

\item Villarroel Enriquez (2025), en su tesis titulada Análisis dinámico de malware mediante algoritmos de detección basados en machine learning, ejecutó una investigación experimental orientada a fortalecer la seguridad en dispositivos móviles a través del análisis de comportamientos dinámicos. La metodología empleada consistió en la comparación de los algoritmos XGBoost, LightGBM y Random Forest, utilizando la frecuencia de llamadas al sistema como métrica principal. Los hallazgos principales destacou que el algoritmo LightGBM superó a los demás con una precisión del 94.1%, demostrando ser una herramienta eficiente para mitigar amenazas en tiempo real. Este antecedente es relevante pues valida el uso de LightGBM en el contexto nacional, aunque su enfoque en malware deja espacio para investigar su desempeño en la detección específica de intrusiones en redes IoT.

\item León et al. (2024), en el trabajo Sistema que alerta la presencia de anomalías en logs de equipos IoT para prevenir el acceso no autorizado, propusieron una solución tecnológica para la detección temprana de vulnerabilidades en infraestructuras conectadas. El método consistió en el desarrollo de un sistema autónomo capaz de analizar grandes volúmenes de logs de auditoría mediante scripts de Python. Los resultados evidenciaron una efectividad variable (65% al 100%) supeditada a la calidad y cantidad de los datos recolectados de los dispositivos. La relevancia de este estudio reside en su aplicación práctica en entornos IoT reales; no obstante, su análisis se limita a registros de eventos, lo que genera un vacío respecto al análisis de flujos de red que la presente investigación cubrirá.

\item Oblitas Mantilla y Escobedo Cárdenas (2024), en Prediction of PM2.5 and PM10 concentrations using XGBoost and LightGBM algorithms: A Case Study in Lima, Peru, realizaron un estudio cuantitativo aplicado a la predicción de variables ambientales complejas. Metodológicamente, los autores compararon el desempeño de XGBoost y LightGBM procesando series temporales masivas obtenidas en el distrito de San Borja. Los resultados demostraron que LightGBM ofrece una velocidad de entrenamiento superior y una mayor precisión en el manejo de datos voluminosos. Aunque el dominio de aplicación es ambiental, este antecedente es fundamental porque proporciona el sustento técnico y la validez metodológica en el Perú para la selección de XGBoost y LightGBM como los pilares del modelo de ensamble propuesto en esta tesis.

\item Curioso Mercado y Gonza Callenova (2025), en Propuesta de modelo de detección de ransomware basado en Aprendizaje Automático, desarrollaron una investigación experimental con el objetivo de identificar ataques de secuestro de datos antes de su ejecución. La metodología se centró en el preprocesamiento de características de red y el entrenamiento de modelos de clasificación supervisada para distinguir tráfico legítimo de patrones de cifrado. Los hallazgos indicaron que la optimización de hiperparámetros permite reducir drásticamente los falsos negativos en amenazas críticas. Este estudio se conecta con la presente investigación al subrayar la importancia del Recall, pero deja pendiente la mejora en la detección de ataques minoritarios, aspecto que se resolverá con el enfoque híbrido planteado.

\item Concha Ramos (2024), en Detección de intrusos en la red basado en red neuronal convolucional y aprendizaje por transferencia, ejecutó un estudio de nivel maestría para analizar el rendimiento de arquitecturas de aprendizaje profundo frente a ataques emergentes. El diseño metodológico empleó el dataset UNSW-NB15 y aplicó técnicas de Transfer Learning sobre Redes Neuronales Convolucionales (CNN) para evitar el reentrenamiento constante. Los resultados arrojaron una precisión del 85.22%, concluyendo que estos modelos son adaptables pero computacionalmente exigentes. La importancia de este antecedente radica en que establece un punto de comparación donde los modelos de ensamble basados en árboles (como el propuesto en este proyecto) suelen ofrecer una mayor exactitud y menor latencia en redes IoT.

\end{enumerate}

\subsection{Antecedentes locales}
\begin{enumerate}
\item 
\end{enumerate}

\section{Bases Teóricas}

\subsection{Internet de las Cosas (IoT) y Ciberseguridad}
El Internet de las Cosas (IoT) se define como una red de objetos físicos interconectados que recolectan e intercambian datos. En este entorno, la seguridad es crítica debido al desbalance de datos en el tráfico de red, donde los ataques maliciosos son significativamente menos frecuentes que el tráfico benigno, dificultando su detección asertiva por parte de los sistemas convencionales.

\subsection{Algoritmos de Gradiente Incrementado (Boosting)}
Los algoritmos de Boosting construyen modelos de forma secuencial, donde cada nuevo modelo intenta corregir de manera iterativa los errores cometidos por los modelos anteriores, optimizando una función de pérdida específica.

\subsection{XGBoost (Extreme Gradient Boosting)}
Es una implementación optimizada del algoritmo de gradiente incrementado que utiliza un término de regulación para controlar el sobreajuste y mejorar la eficiencia computacional. Su función objetivo a optimizar se define como:

\begin{equation}
\mathcal{L}(\phi) = \sum_{i} l(\hat{y}_i, y_i) + \sum_{k} \Omega(f_k)
\end{equation}

Donde $l$ representa la función de pérdida que mide la diferencia entre la predicción $\hat{y}_i$ y el valor real $y_i$, mientras que $\Omega$ es el término de regularización que penaliza la complejidad del modelo para evitar el overfitting.

\subsection{LightGBM (Light Gradient Boosting Machine)}
Este algoritmo se diferencia por su técnica de crecimiento de árboles por hojas (\textit{leaf-wise}) en lugar de niveles (\textit{level-wise}). Este enfoque permite reducir la pérdida de forma más agresiva, mejorando la velocidad de entrenamiento y disminuyendo el uso de memoria en el procesamiento de grandes volúmenes de datos.

\subsection{Ensamble por Stacking (Apilamiento)}
El \textit{Stacking} es una técnica de aprendizaje de ensamble que combina las predicciones de múltiples modelos base (en este estudio, XGBoost y LightGBM) mediante un meta-clasificador superior. El meta-modelo aprende la configuración óptima para combinar las salidas de los modelos base, maximizando así la capacidad de generalización y la robustez del sistema final.

\subsection{Técnicas de Tratamiento de Desbalance de Clases}
\subsection{Borderline-SMOTE}
Es una variante avanzada del algoritmo SMOTE (\textit{Synthetic Minority Over-sampling Technique}) que genera ejemplos sintéticos exclusivamente en la ``frontera'' de decisión entre la clase minoritaria y la mayoritaria. Esta técnica es fundamental cuando el riesgo de clasificación errónea es elevado debido a la proximidad de ejemplos de diferentes clases en el espacio de características.

\subsection{Métricas de Evaluación de Clasificación}
Para evaluar la eficacia del modelo en la detección de ataques minoritarios, se utilizan métricas derivadas de la matriz de confusión, permitiendo una interpretación objetiva del rendimiento:

\begin{itemize}
    \item \textbf{Accuracy (Exactitud):} Proporción de predicciones correctas sobre el total de casos evaluados.
    \begin{equation}
    Accuracy = \frac{TP + TN}{TP + TN + FP + FN}
    \end{equation}

    \item \textbf{Recall (Sensibilidad):} Capacidad del modelo para identificar correctamente los casos positivos reales (ataques). Es la métrica crítica en esta investigación para reducir los falsos negativos.
    \begin{equation}
    Recall = \frac{TP}{TP + FN}
    \end{equation}

    \item \textbf{F1-Score:} Media armónica entre la Precisión y el Recall, proporcionando una medida equilibrada del rendimiento en presencia de clases desbalanceadas.
    \begin{equation}
    F1\text{-}Score = 2 \times \frac{Precision \times Recall}{Precision + Recall}
    \end{equation}
\end{itemize}

\section{Definición de Términos Básicos}

\begin{description}
    \item[\textbf{Ataque Minoritario:}] Actividad maliciosa cuya frecuencia de aparición dentro de un conjunto de datos es significativamente baja en comparación con el volumen del tráfico normal o benigno.
    \item[\textbf{Dataset CICIoT2023:}] Conjunto de datos estandarizado y actualizado, diseñado por el Canadian Institute for Cybersecurity, que contiene perfiles de ataques reales y actualizados en ecosistemas de dispositivos IoT.
    \item[\textbf{Hiperparámetros:}] Configuraciones externas al modelo que se definen previo al proceso de entrenamiento (como la tasa de aprendizaje o profundidad del árbol) para optimizar su rendimiento.
    \item[\textbf{Recall:}] Indicador métrico que cuantifica la habilidad de un clasificador para recuperar y etiquetar correctamente todos los casos positivos existentes en la muestra.
    \item[\textbf{Stacking:}] Arquitectura de aprendizaje jerárquico donde un meta-modelo se entrena utilizando las predicciones de diversos modelos base para reducir el error de sesgo y varianza.
    \item[\textbf{SMOTE:}] Algoritmo de sobremuestreo que incrementa artificialmente la representatividad de la clase minoritaria mediante la síntesis de nuevos datos basados en los vecinos más cercanos.
    \item[\textbf{Verdaderos Positivos (TP):}] Cantidad de instancias en las que el modelo predijo correctamente la existencia de un ataque o anomalía.
\end{description}